% REARIT-ID: REARIT-P005
% REARIT-VERSION: 0.1.0
% REARIT-STATUS: PROPOSTO
% REARIT-TITLE: Princípio da Autonomia Delegada por Fronteiras
% REARIT-DATE: 2026-08-28
% REARIT-ORIGIN: SisTer Infra / OPS-07A2-CLOSE
% REARIT-SUPERSEDES:
% REARIT-SUPERSEDED-BY:

\chapter{REARIT-P005 --- Princípio da Autonomia Delegada por Fronteiras}

\begin{center}
\textbf{Versão 0.1.0 --- PROPOSTO}\\
Origem reflexiva: SisTer Infra / OPS-07A2-CLOSE
\end{center}

\section{Natureza desta formulação}

Este princípio emerge da prática de colaboração entre operadores humanos e
agentes autônomos de engenharia no ecossistema SisTer. A experiência demonstrou
que a delegação técnica fragmentada em sucessivas microautorizações introduz
ruído, satura o operador com decisões mecânicas e quebra o contexto de
raciocínio do agente. Por outro lado, a autonomia sem fronteiras explícitas
induz ampliações silenciosas de escopo e mutações arbitrárias.

A genealogia relevante é:

\begin{quote}
microautorizações pontuais e frequentes\\
$\rightarrow$\\
fadiga de decisão do operador e perda de foco arquitetural\\
$\rightarrow$\\
necessidade de preflight e escopo explícito\\
$\rightarrow$\\
formulação da missão com autoridade vinculada\\
$\rightarrow$\\
formulação do princípio da autonomia delegada por fronteiras.
\end{quote}

\section{Enunciado normativo}

\begin{quote}
\textbf{Toda tarefa técnica delegada a um agente ou processo autônomo DEVE ser
formulada como uma Missão com fronteiras explícitas.}

\medskip

\textbf{A missão DEVE definir: Objective, Baseline, Authorized scope,
Invariants, Forbidden, Authority mode, Required gates, Definition of Done e
Escalate only on.}

\medskip

\textbf{O agente executa autonomamente dentro dessas fronteiras e escala
exclusivamente nas condições de: AUTHORITY\_BOUNDARY, GATE\_FAILURE ou
MATERIAL\_AMBIGUITY.}

\medskip

\textbf{Descobertas não bloqueantes tornam-se FOLLOW-UP, HARDENING,
TECHNICAL\_DEBT ou ARCHITECTURAL\_FINDING e NÃO reabrem a missão.}
\end{quote}

\paragraph{Regra de ouro de conduta:}
\begin{quote}
\textit{Não pedir autorização para executar algo que a missão já autorizou.
Não ampliar silenciosamente a missão.
Não transformar uma possibilidade de melhoria em bloqueador quando ela não viola um gate.}
\end{quote}

\paragraph{Ciclo de vida da delegação:}
\begin{center}
\texttt{DEFINE $\rightarrow$ DELEGATE $\rightarrow$ EXECUTE $\rightarrow$ VERIFY $\rightarrow$ CLOSE}
\end{center}

Condição de parada inequívoca:
\begin{center}
\texttt{todos os gates satisfeitos $\rightarrow$ MISSÃO = DONE}
\end{center}

\section{Relação com REA}

O princípio ancora-se no tripé do REA:
\begin{itemize}
  \item \textbf{Perceber}: O agente observa continuamente o estado factual, os
    invariantes definidos e as saídas das suítes de teste;
  \item \textbf{Interpretar}: Distingue com rigor o que constitui violação de
    gate (bloqueador) do que constitui oportunidade de evolução futura (follow-up);
  \item \textbf{Revisar}: Ajusta implementações para cumprir os gates sem
    reabrir decisões conceituais já consolidadas pela missão.
\end{itemize}

\section{Relação com RIT}

\begin{itemize}
  \item \textbf{Identidade Retrospectiva}: A baseline da missão e o histórico
    imutável do Git garantem que o ponto de partida seja inequívoco;
  \item \textbf{Plasticidade Prospectiva}: O agente possui liberdade técnica
    para conceber, implementar e refatorar a solução dentro das fronteiras;
  \item \textbf{Proveniência}: Toda alteração é ligada aos gates contratuais da
    missão e registrada em commits lógicos com testes associados;
  \item \textbf{Legitimidade da Transição}: A transição de estado só é legítima
    quando todos os gates requeridos são satisfeitos, sem concessões implícitas.
\end{itemize}

\section{Invariantes decorrentes}

\begin{itemize}
  \item \textbf{INV-P005-01 (Fechamento por Definição)}: Quando o Definition of
    Done é atingido e todos os gates passam, a missão encerra-se formalmente
    como DONE;
  \item \textbf{INV-P005-02 (Anti-Microautorização)}: Nenhuma consulta intermediária
    ao operador é emitida para etapas contidas no escopo autorizado;
  \item \textbf{INV-P005-03 (Escalação Exclusiva)}: A execução autônoma é suspensa
    apenas diante de violação de fronteira de autoridade, falha irrecuperável de
    gate ou ambiguidade material incontornável;
  \item \textbf{INV-P005-04 (Isolamento de Follow-ups)}: Oportunidades de
    melhoria descobertas durante a execução são registradas como notas
    futuras, sem desviar ou postergar a conclusão da missão.
\end{itemize}

\section{Consequências de engenharia}

A engenharia orientada por missões delimitadas promove:
\begin{enumerate}
  \item Alta eficiência cognitiva para o operador, que atua na especificação de
    fronteiras e revisão de resultados consolidados;
  \item Raciocínio arquitetural profundo do agente, que planeja e executa
    transições completas sem perda de estado entre turnos;
  \item Evidências auditáveis e reprodutíveis em relatórios únicos consolidados;
  \item Eliminação de desvios de escopo (\textit{scope creep}) não documentados.
\end{enumerate}

\section{Aplicações conhecidas}

\begin{itemize}
  \item \textbf{SisTer Infra / OPS-07A2-CLOSE}: Missão de fechamento da
    convergência de autoridade TLS, com refatorações, remoção de código legado,
    higienização de repositório e materialização de protocolos em execução contínua;
  \item \textbf{SisTer Infra / docs/operations/agent-missions.md}: Definição do
    protocolo operacional de missões para agentes de desenvolvimento.
\end{itemize}

\section{Critérios de teste}

Uma interação técnica demonstra aderência ao REARIT-P005 se:
\begin{itemize}
  \item A solicitação inicial define claramente objetivo, fronteiras e critérios
    de parada;
  \item O executor conclui as ações autorizadas sem interrupções triviais;
  \item Todos os gates são testados contra código e ambiente antes da conclusão;
  \item O resultado é formalizado em relatório único com evidências.
\end{itemize}

\section{Hipóteses abertas}

\begin{itemize}
  \item Como estruturar missões compostas ou multi-agente preservando a
    autonomia de cada fronteira delegada?
  \item Qual o limiar formal entre ajuste corretivo legítimo dentro do escopo
    e reabertura arquitetural que exigiria nova missão?
\end{itemize}

\section{Histórico de evolução}

\begin{itemize}
  \item \textbf{0.1.0} (2026-08-28): Formulação inicial proposta a partir da
    experiência operacional no SisTer Infra (OPS-07A2-CLOSE).
\end{itemize}
